% TOP_PROBLEM: {"problemId": "problem.shelah-s-categoricity-conjecture-for-l-omega1-omega", "canonicalTitle": "Shelah's Categoricity Conjecture for L_omega1,omega", "releaseRank": 370, "primaryDomain": "logic_foundations_set_theory", "primaryDomainLabel": "Logic, foundations & set theory", "displayStatus": "open_with_solved_subcases", "releaseStatus": "open_with_solved_subcases"}
% !TeX program = lualatex
% ATTEMPT_STATUS: unresolved
\documentclass[11pt]{article}
\usepackage[margin=25mm]{geometry}
\usepackage{fontspec}
\setmainfont{Latin Modern Roman}
\usepackage{amsmath,amssymb,mathtools}
\usepackage{unicode-math}
\setmathfont{Latin Modern Math}
\usepackage{xurl,hyperref}
\hypersetup{colorlinks=true,urlcolor=blue}
\setcounter{secnumdepth}{0}
\setlength{\parindent}{0pt}
\setlength{\parskip}{5pt}
\setlength{\emergencystretch}{3em}
\begin{document}
\section*{\texorpdfstring{Shelah's Categoricity Conjecture for L\_omega1,omega}{Shelah's Categoricity Conjecture for L\_omega1,omega}}\label{370}
\subsection{Definitions and mathematical statement}
The logic $L_{\omega_1,\omega}$ permits countable conjunctions and disjunctions but only finite strings of quantifiers. A sentence $\psi$ is categorical in $\lambda$ if it has exactly one model of cardinality $\lambda$ up to isomorphism. \emph{Source notation correction:} Conjecture 1.1 of Shelah--Vasey prints the threshold $\beth_{\omega_1}$, not the undefined ``delta\_omega1'' of the JSON transcription. Define $\beth_0=\aleph_0$, $\beth_{\alpha+1}=2^{\beth_\alpha}$, and $\beth_\lambda=\sup_{\alpha<\lambda}\beth_\alpha$ at limit ordinals. The cited strong conjecture is
\[
 \bigl[\exists\lambda\ge\beth_{\omega_1}:\psi\text{ categorical in }\lambda\bigr]
 \Longrightarrow
 \bigl[\forall\mu\ge\beth_{\omega_1}:\psi\text{ categorical in }\mu\bigr].
\]
The correction is based on the source's displayed statement. Allowing an unspecified larger eventual threshold would weaken this question.
\par\smallskip\textit{Formulation and concept references:} \href{https://shelah.logic.at/files/241811/10.4171-jems-1477.pdf}{[S1]}.

\subsection{Short English statement}
If an infinitary sentence has a unique model at one cardinal above the specified Hanf threshold, must it have a unique model at every cardinal above that threshold?
\subsection{Research attempt}
\paragraph{The exact threshold and the source audit (27 September 2026).}
The target is the strong version with the fixed threshold
$\beth_{\omega_1}$, not eventual categoricity above an unspecified larger
cardinal. The supplied Shelah--Vasey paper [S1], Conjecture 1.1 and the
following discussion, were inspected. Its Corollary 14.7 proves an
eventual statement under a strongly compact cardinal hypothesis and
with a threshold depending on that cardinal. Those extra assumptions
and the different threshold cannot be erased to solve the selected
ZFC statement.

We use the usual countable vocabulary convention for the infinitary
sentence. The countably many symbols and subformulas actually appearing
in a sentence are the relevant syntactic data; arbitrary expansions
by unrelated new symbols are not part of the categoricity assertion.
Categoricity includes existence: zero models in a cardinal is not
categoricity there.

The attempt develops three checks. First, a complete vector-space
family illustrates a successful dimension argument. Second, explicit
infinitary sentences show what countable categoricity and compactness
cannot give. Third, a fragment-hull construction and a precise
back-and-forth lemma isolate the missing extension property. None
of the counterexamples below satisfies the large-cardinal premise of
the conjecture; they test proposed shortcuts, not the conjecture itself.

\paragraph{A fully proved categorical family in every infinite cardinal.}
Fix a finite or countably infinite field $F$. Use symbols for addition,
zero, and multiplication by each scalar in $F$. The vector-space axioms
form a countable set of first-order sentences, so their conjunction is
an $L_{\omega_1,\omega}$ sentence. Conjoin, for each positive integer
$n$, the assertion that there are $n$ linearly independent vectors.
For a countable field, independence of a fixed finite tuple is expressed
by the countable conjunction
\[
 \bigwedge_{(a_1,\ldots,a_n)\in F^n\setminus\{0\}}
                     a_1x_1+\cdots+a_nx_n\ne0.
\]
Only $n$ existential quantifiers occur in that assertion; the full
sentence uses a countable conjunction of these finite strings, not
an illegal infinite quantifier string. Denote the resulting sentence
by $\psi_F$.

Its models are exactly the infinite-dimensional $F$-vector spaces.
Every vector space has a basis by the usual Zorn argument: order linearly
independent subsets by inclusion, take a maximal one, and note that a
vector outside its span could be adjoined. A bijection between bases
extends uniquely by finite linear combinations to a vector-space
isomorphism. Thus dimension determines the isomorphism type.

If the basis has infinite cardinality $\delta$ and $F$ is countable,
the set of finite linear combinations has cardinality
$\max(\delta,\aleph_0)=\delta$. The upper bound follows by a countable
union over finite tuple lengths, and the lower bound follows from
containing the basis. Consequently a model of infinite cardinality
$\lambda$ has dimension $\lambda$, and the direct sum
$F^{(\lambda)}$ supplies one. Therefore
\[
 I(\psi_F,\lambda)=1\qquad\text{for every infinite }\lambda.
 \tag{370.1}
\]
Here $I$ counts isomorphism types of models of the indicated size.
This verifies the selected conclusion for an explicit class, including
all cardinals above the requested threshold.

The argument has a genuine invariant, basis cardinality, together with
an extension theorem for maps of bases. Merely calling an arbitrary
model's cardinality its ``dimension'' would not reproduce either step.

\paragraph{Why finite dimension over a countable field needs care.}
If $F$ is countably infinite, a nonzero finite-dimensional vector space
is already countably infinite as a set. Omitting the infinite-dimensional
conjunct would therefore give several nonisomorphic countable models,
although the uncountable models are still classified by their cardinality.
The sentence above includes the conjunct explicitly so that (370.1)
is valid also at $\aleph_0$.

This small example matters when using cardinal arithmetic in a general
argument: the formula $|V|=\dim_FV$ was proved only for infinite
basis cardinality. A countable language also does not force all
countable models of a theory to be isomorphic.

\paragraph{A sentence categorical only in the countable cardinal.}
Take countably many constant symbols $c_0,c_1,\ldots$, with no other
nonlogical symbols, and form
\[
 \chi=\left(\bigwedge_{i<j}c_i\ne c_j\right)
       \wedge\forall x\left(\bigvee_{n<\omega}x=c_n\right).
 \tag{370.2}
\]
Every model has exactly the elements named by the distinct constants,
so it is countably infinite. Between two such models the unique map
sending $c_n$ to $c_n$ is a bijection and an isomorphism. Thus $\chi$
is categorical in $\aleph_0$ and has no uncountable models.

This does not challenge the conjecture, whose premise begins at
$\beth_{\omega_1}$. It directly refutes an argument that transfers
categoricity from an arbitrary infinite cardinal merely by adding
unnamed elements. Formula (370.2) prohibits those elements in one
countable disjunction.

It also demonstrates failure of compactness in the required logic.
Add a new constant $d$ and the sentences $d\ne c_n$ for all $n$.
Every finite part of this theory has a model: retain $\chi$ and
interpret $d$ as a constant not among the finitely many forbidden
ones. The entire theory has no model. First-order compactness cannot
be applied to this collection as if its infinitary disjunction were
a finitary formula.

\paragraph{Even existence in all cardinals does not repair a countable argument.}
The first-order theory of dense linear orders without endpoints is
categorical in $\aleph_0$. A direct proof is the classical alternating
extension construction: enumerate the two countable orders, and extend
a finite order-preserving partial bijection alternately to include the
next unused element of each. Density and the absence of endpoints
supply an image in the required finite interval. The union is an
isomorphism.

For every uncountable cardinal $\lambda$, that same theory has
nonisomorphic models of size $\lambda$. Here is an explicit pair.
Let $J=\lambda\times\mathbb Q$ with lexicographic order. It is dense
and has no endpoints: within a block use density of $\mathbb Q$,
and between different blocks a larger rational in the earlier block
lies between the two specified points. Put
\[
 L_0=\mathbb Z\times J,\qquad L_1=\omega_1\times J
\]
with lexicographic orders. Both are dense without endpoints and both
have cardinality $\lambda$.

The cofinality of $L_0$ is $\aleph_0$: choose one point in each
positive integer block to obtain a countable cofinal sequence, and
no finite set is cofinal in an order with no last point. The cofinality
of $L_1$ is $\aleph_1$. Any countable set has its first coordinates
bounded in $\omega_1$, so is not cofinal. Choosing one point in each
block is cofinal and has size $\aleph_1$. Cofinality is preserved by
order isomorphisms, so these models are not isomorphic.

The finite-partial-map proof for countable orders failed because at
uncountable stages one must realize cuts determined by potentially
infinite earlier domains. Density only promises points between finitely
many bounds. This is a concrete model of the extension problem a
transfinite categoricity proof must solve; increasing the length of
the enumeration does not automatically upgrade the available extension
property.

\paragraph{A fragment version of downward Löwenheim--Skolem can be proved directly.}
Let $\mathcal F$ be a countable collection of
$L_{\omega_1,\omega}$ formulas containing the sentence $\psi$ and
closed under subformulas and the finitely many-variable operations
needed for the following witness argument. Such a countable fragment
can be chosen for one sentence by recursively adjoining its subformulas
and closing under countably many syntactic operations. Let
$N\models\psi$ and $A\subseteq N$.

Build sets $A_0\subseteq A_1\subseteq\cdots$ starting with $A$
and the interpretations of constants. At each step close under all
function symbols, and for each existential formula
$\exists x\,\varphi(x,\bar a)\in\mathcal F$ with parameters in
the current set that is true in $N$, choose one witness in $N$ and
adjoin it. There are at most $|A_i|+\aleph_0$ choices, since the
fragment is countable and each parameter tuple is finite. Put
$M=\bigcup_{i<\omega}A_i$ with the induced structure. Then
\[
 |M|\le |A|+\aleph_0.
 \tag{370.3}
\]

Induction on formulas in the fragment shows that $M$ and $N$ agree
on their truth values for parameters from $M$. Atomic formulas agree
because it is a substructure. Negation and finite Boolean operations
preserve agreement. For a countable conjunction, agreement holds for
each conjunct by induction, so holds for their conjunction; the same
applies to disjunction. For an existential formula true in $N$, its
finitely many parameters appear together at some finite construction
stage, and the next stage supplies a witness in $M$. The converse
follows from agreement for the witnessing subformula. Thus
$M\preccurlyeq_{\mathcal F}N$, and in particular $M\models\psi$.

For an infinite $\mu\le|N|$, starting with a set $A$ of size $\mu$
therefore gives a model of $\psi$ of size exactly $\mu$. This yields
\emph{existence} below a known model. It neither constructs larger
models nor shows that the models of a given smaller size are isomorphic.
The submodel relation is elementary for the specified fragment, not
for every formula in the entire infinitary language.

\paragraph{Unions require a compatible elementary-chain hypothesis.}
The same fragment induction proves a useful chain statement. If
$M_i\preccurlyeq_{\mathcal F}M_j$ for $i<j$ in an increasing chain,
then each $M_i$ is $\mathcal F$-elementary in the union. In the
existential step, a witness in the union belongs to some later stage,
and elementarity transfers the existential assertion back to the stage
containing its parameters. A witness at that stage also works in the
union. The Boolean and countable-conjunction steps are pointwise.

But arbitrary embeddings between models of $\psi$ need not be
$\mathcal F$-elementary. Nor does a sequence of isomorphism types
come equipped with chosen compatible embeddings. Building a directed
union without specifying a relation that guarantees preservation of
$\psi$ is therefore not a proof of existence, much less categoricity,
at a limit cardinal. This is one reason the model-theoretic source
works with controlled classes and embeddings rather than only raw
cardinality counts.

\paragraph{A precise sufficient extension principle for uniqueness.}
\textbf{Lemma 370.1.} Let $M,N$ be structures of the same infinite
cardinality $\kappa$. Suppose there is a nonempty family $\mathcal I$
of partial isomorphisms from $M$ to $N$, each of domain and range size
less than $\kappa$, with the following properties. The family is closed under restrictions.
Any member can be
extended in $\mathcal I$ to include any prescribed element of $M$
in its domain and any prescribed element of $N$ in its range. Unions
of increasing chains from $\mathcal I$ whose union has size less than
$\kappa$ also belong to $\mathcal I$. Then $M\cong N$.

\emph{Proof.} Enumerate each universe in length $\kappa$. Recursively
build an increasing sequence of partial maps, at successor stages
including the next listed domain and range elements by the two extension
properties. At a limit stage below $\kappa$, take the union.
Start with the empty map, available by restriction closure. After each
extension restrict back to the old domain and the at most two new
elements needed for the requests. At stage $\alpha$ the domain has
size at most $2|\alpha|$, hence is finite when $\alpha$ is finite
and is less than $\kappa$ for every $\alpha<\kappa$. The union
hypothesis therefore applies, even for singular $\kappa$.

At the final stage the union is a total bijection. Every atomic relation
and function instance involves finitely many elements, hence is covered
by one earlier partial isomorphism. The union is an isomorphism.
\hfill$\square$

The restriction hypothesis is essential to the displayed size control.
An extension between whole submodels may add arbitrarily many elements
below $\kappa$; an increasing sequence of such sizes can become
cofinal in a singular $\kappa$ too early. The lemma is a sufficient
criterion for the precisely specified partial-map family, not a claim
that every family of submodel embeddings satisfies its hypotheses.

\paragraph{Where applying this lemma to the conjecture stops.}
The vector-space example supplies such maps by mapping small independent
sets and extending maps of bases. The countable dense-order proof
supplies finite maps, enough for $\kappa=\aleph_0$. A general sentence
categorical at one huge cardinal does not visibly provide any such
family for arbitrary models at each other huge cardinal.

An isomorphism between two large models is also not automatically an
isomorphism fixing a specified common submodel pointwise. The latter
requires homogeneity over that submodel, or an amalgamation/uniqueness
argument controlling the embeddings. Existence of an isomorphism type
and compatibility of diagrams are distinct assertions. The multidimensional
diagrams in [S1] address precisely stronger structural information; their
hypotheses cannot be replaced with an informal appeal to cardinality.

Even an eventual theorem beginning above a larger cardinal $\theta$
would leave the interval $[\beth_{\omega_1},\theta)$ untreated if
$\theta>\beth_{\omega_1}$. The downward fragment construction gives
models in that interval but no uniqueness theorem there. Thus it cannot
repair the difference between the source's conditional eventual result
and the sharp threshold required in this record.

\paragraph{Outcome and first missing step.}
The arguments were checked against explicit vector-space bases, the
countable-only sentence (370.2), and the cofinality distinction between
the two dense orders. These are exact logical examples. Finite enumeration
of small structures cannot test categoricity at the beth threshold and
is not presented as verification of such a claim.

\textbf{Outcome: UNRESOLVED.} The complete positive family has a genuine
dimension-and-extension theorem, and the fragment construction supplies
controlled downward existence and elementary unions. For an arbitrary
sentence, the missing step is to derive sufficient structural extension
and uniqueness properties from the one-cardinal categoricity premise,
with the exact threshold and without additional large-cardinal axioms.
None of compactness, countable back-and-forth, or a larger eventual
threshold supplies that step as written.

\subsection{Sources}
\begin{itemize}
\item[S1]Shelah and Vasey, Categoricity and multidimensional diagrams, JEMS, 2024.
\url{https://shelah.logic.at/files/241811/10.4171-jems-1477.pdf}.
\textit{Formulation source. Consulted 24 September 2026: Conjecture 1.1, printed page 2302: the threshold is beth\_\{omega\_1\}, not delta\_\{omega\_1\}..}
\end{itemize}
\end{document}
